\chapter{Réalisation de la plateforme}

\section{Organisation du projet}

Le dépôt est organisé par service. Chaque dossier contient sa configuration,
son code et un fichier README.

\begin{lstlisting}[language=bash,caption={Organisation des principaux composants}]
.
|-- broker/        # Eclipse Mosquitto
|-- dashboards/    # provisioning et tableaux Grafana
|-- simulator/     # modele Python de l'imprimante
|-- ingestion/     # image et flux Node-RED
|-- notifier/      # diffusion des alertes
|-- docker-compose.yml # services, volumes et initialisation InfluxDB
`-- Makefile
\end{lstlisting}

Les services de base démarrent sans profil. Le simulateur, Node-RED et le
notifier utilisent des profils Compose. La commande \texttt{make demo}
construit et lance l'ensemble.

\section{Simulateur de l'imprimante CIJ}

Le simulateur est écrit en Python. Il contient un état interne, deux
consommables et un générateur de commandes. Son cycle de fonctionnement suit
les états \texttt{shutdown}, \texttt{warmup}, \texttt{running},
\texttt{stopped} et \texttt{fault}. Une loi de Poisson simplifiée déclenche des
défauts à une fréquence configurable.

Le modèle reprend les données documentées pour la 9450c : cartouches externes
de 800~mL, réservoirs internes de 1,075~L, état du jet, état d'impression,
vitesse moteur, pression, viscosité, températures, niveaux et autonomie
d'encre. Les codes natifs sont publiés en parallèle de l'état
fonctionnel utilisé par le tableau de bord. Par exemple, le code jet
\texttt{00h} représente l'arrêt, \texttt{01h} le démarrage et \texttt{07h} le
fonctionnement normal ; l'état d'impression distingue pause, impression
opérationnelle et imprimante non prête.

\begin{table}[H]
\centering
\caption{Comportements reproduits par le simulateur}
\begin{tabularx}{\textwidth}{p{0.34\textwidth}Y}
\toprule
Comportement & Utilité pour les essais \\
\midrule
Périodes de commande et périodes vides &
Vérifier la disponibilité liée à la demande. \\
Consommation d'encre en production &
Tester les tendances et la prévision d'épuisement. \\
Évaporation du solvant avec le jet actif &
Éviter une prévision basée uniquement sur le nombre de pièces. \\
Remplissage avec saut de niveau &
Vérifier le changement de pente et la gestion des lots. \\
Dérive de viscosité &
Déclencher une alerte avant une dégradation visible. \\
Paramètres natifs de la 9450c &
Vérifier le mapping futur des trames constructeur vers Litmus Edge. \\
Défauts avec code et durée &
Alimenter l'historique, le temps d'arrêt et le Pareto. \\
Silence temporaire du producteur &
Tester l'alerte d'absence de données. \\
\bottomrule
\end{tabularx}
\end{table}

Chaque message contient un UUID dans \texttt{message\_id}. Le Last Will MQTT
annonce aussi un arrêt si le client disparaît sans fermeture propre.

\section{Ingestion avec Node-RED}

\begin{figure}[H]
  \centering
  \includegraphics[width=\textwidth]{flux-node-red.pdf}
\caption{Chaîne d'ingestion mise en œuvre avec Node-RED}
  \label{fig:nodered}
\end{figure}

Node-RED contient trois flux de données et un flux de supervision :

\begin{enumerate}
  \item \textbf{Télémétrie} reçoit les valeurs continues, contrôle le payload
        et construit un point \texttt{printer\_telemetry}.
  \item \textbf{Événements} traite les changements d'état, défauts,
        remplissages et annonces de démarrage.
  \item \textbf{Marquages} enregistre le DMC et le résultat de contrôle.
  \item \textbf{Supervision} sonde InfluxDB, expose \texttt{/health} et
        centralise les erreurs.
\end{enumerate}

\begin{figure}[H]
  \centering
  \includegraphics[width=\textwidth]{node-red-telemetrie.png}
  \caption{Flux Node-RED de télémétrie : MQTT, validation et écriture InfluxDB}
  \label{fig:nodered-telemetrie-reel}
\end{figure}

Node-RED utilise un nœud HTTP natif pour appeler directement l'API d'écriture
v2 d'InfluxDB. Aucun module de base de données supplémentaire n'est nécessaire.
Le token, le nom
de l'organisation et le bucket proviennent des variables d'environnement ; ils
ne sont pas inscrits dans \texttt{flows.json}. Le nœud de préparation échappe
les tags et les chaînes avant de produire le Line Protocol :

\begin{lstlisting}[language={},caption={Exemple de Line Protocol produit par Node-RED}]
printer_telemetry,printer_id=CIJ_Printer_L1,line_id=L1 \
ink_level=74.1,solvent_level=57.3,ink_viscosity=4.19 1785445200000
\end{lstlisting}

L'appel HTTP utilise une précision en millisecondes. Il cible l'organisation
\nolinkurl{saint-gobain} et le bucket \nolinkurl{printer_monitoring}. Un message
invalide est publié sur le topic de rejet avec son topic d'origine et un motif
compréhensible par l'exploitant.

\begin{figure}[H]
  \centering
  \includegraphics[width=\textwidth]{node-red-supervision.png}
  \caption{Flux Node-RED de supervision et de traitement des rejets}
  \label{fig:nodered-supervision-reel}
\end{figure}

\section{Écriture idempotente et reprise}

InfluxDB identifie un point par sa mesure, ses tags et son horodatage. Le
simulateur conserve l'horodatage d'origine lors d'une retransmission. Node-RED
reconstruit donc la même clé et InfluxDB met à jour le point existant au lieu de
dupliquer l'observation~\cite{influxduplicate}. Le \texttt{message\_id} reste
présent dans les événements pour faciliter le diagnostic, sans être utilisé
comme tag sur la télémétrie afin de limiter la cardinalité.

Une écriture InfluxDB n'offre pas de transaction entre plusieurs mesures. Le
flux traite chaque message de manière indépendante et conserve sa clé lors
d'un nouvel essai. Cette limite est acceptable ici : les messages MQTT sont
autonomes et le rejouage d'un point ne change pas le résultat final.

\section{Requêtes Flux}

Les calculs partagés sont écrits en Flux dans les tableaux de bord et les
règles d'alerte. Ils couvrent notamment :

\begin{itemize}
  \item la dernière valeur de chaque champ de télémétrie ;
  \item la lecture de la classification calculée par Node-RED à partir de la
        demande et de l'état machine ;
  \item la lecture des prévisions d'épuisement et de péremption ajoutées au
        point de télémétrie ;
  \item la disponibilité quotidienne, calculée uniquement pendant la demande ;
  \item le Pareto des défauts et le rendement au premier passage.
\end{itemize}

Les scripts Flux restent dans les fichiers JSON provisionnés de Grafana. Ils
sont contrôlés par la commande \texttt{make verify}, qui exécute les requêtes
contre le bucket de démonstration.

\section{Conteneurisation}

Docker Compose démarre six conteneurs : Mosquitto, InfluxDB, Grafana, le
simulateur, Node-RED et le notifier. Les services attendent la santé de leurs
dépendances. Les données du courtier, d'InfluxDB et de Grafana sont placées
dans des volumes.

Les mots de passe sont générés dans un fichier \texttt{.env} exclu du dépôt.
Le déploiement expose uniquement les interfaces nécessaires aux utilisateurs ;
les services de transport et de stockage restent internes à la plateforme.

\section{Notification des alertes}

Grafana envoie toutes les alertes vers un webhook unique. Le notifier transforme
le payload technique en une phrase courte, puis le transmet aux canaux activés.
Les adaptateurs disponibles sont la console, Telegram, le courrier électronique,
Microsoft Teams et une passerelle WhatsApp de démonstration.

L'adaptateur \texttt{open-wa} n'est pas retenu pour la production. Il pilote
WhatsApp Web, dépend d'une session de téléphone et peut cesser de fonctionner
après une mise à jour. Pour une utilisation industrielle, Teams, l'e-mail ou
l'API WhatsApp Business officielle sont plus faciles à maintenir.
